% Copyright (c) 2023-2026 Oriol Cayon, Delft University of Technology
% SPDX-License-Identifier: Apache-2.0
\documentclass[11pt,a4paper]{article}

\usepackage[margin=2.4cm]{geometry}
\usepackage{graphicx}
\usepackage{amsmath,amssymb}
\usepackage{booktabs}
\usepackage{longtable}
\usepackage{xcolor}
\usepackage{caption}
\usepackage{float}
\usepackage{enumitem}
\usepackage{fancyhdr}
\usepackage[hidelinks]{hyperref}
\usepackage{siunitx}
\sisetup{range-phrase={--},range-units=single}

\graphicspath{{../results/demo/}{../results/validation/}{../results/hanging/figures/}}

\definecolor{accent}{HTML}{D55E00}
\definecolor{muted}{HTML}{555555}

\captionsetup{font=small,labelfont=bf}
\setlist{itemsep=2pt,parsep=0pt,topsep=3pt}

\pagestyle{fancy}
\fancyhf{}
\fancyhead[L]{\small\textsc{Billow}}
\fancyhead[R]{\small\thepage}
\renewcommand{\headrulewidth}{0.3pt}

\newcommand{\code}[1]{\texttt{\small #1}}

\title{%
  \vspace{-1.5cm}
  {\Huge\textbf{Billow}}\\[6pt]
  {\large A minimum-energy structural model for soft kites}\\[10pt]
  {\normalsize Timoshenko beams, inflatable tubes, wrinkling membranes,\\
   cables and pulleys --- one potential energy, one NLP}
}
\author{Oriol Cayon \\ \small Delft University of Technology}
\date{\today}

\newcommand{\billowmean}{15.7}
\newcommand{\kfemmean}{12.4}
\newcommand{\billowmeansix}{11.5}
\newcommand{\kfemmeansix}{7.8}
\newcommand{\billowtime}{14}
\newcommand{\billowspanone}{7.276}
\newcommand{\storedmean}{15.1}
\newcommand{\storedspanone}{7.212}
\newcommand{\kfemtime}{552}

\begin{document}
\maketitle
\thispagestyle{fancy}

\begin{abstract}
\noindent
Billow is a standalone structural library for soft-wing airborne wind energy
systems. It assembles cables, frictionless pulleys, geometrically exact
Timoshenko beams, inflatable tubes and wrinkling membrane fabric into a single
total potential energy, and finds static equilibrium by minimising it with
IPOPT. Stationarity of the energy \emph{is} the equilibrium equation, so a
converged solve satisfies the nodal force balance to solver tolerance rather
than to whatever residual a relaxation scheme happens to leave behind.

This document specifies the formulation, the element library, the validation
against benchmarks external to the code, a set of demonstration cases, and a
measured comparison against \code{kite\_fem}, the existing FEM path. Headline
results: exact-to-$10^{-4}$ agreement with the Euler elastica, second-order
convergence to the exact roll-up circle, agreement with published values for
the Bathe--Bolourchi $45^\circ$ bend, reproduction of the ASKITE inflatable-tube
fits to $0.06\%$ in bending, and a full-kite structural solve in $5.9$\,s cold
and $0.58$\,s warm at 1893 degrees of freedom.

Billow currently lives as \code{awetrim.structural} and imports only NumPy and
CasADi. It is intended to become its own repository.
\end{abstract}

\tableofcontents
\newpage

%======================================================================
\section{Scope and design intent}

\subsection{What Billow is}

A static structural solver for the kinds of structures soft kites are made of:
thin lines that cannot push, ropes that run over sheaves, inflated tubes that
bend until they collapse, and fabric that carries tension but wrinkles instead
of carrying compression. Given a geometry, material properties and a set of
nodal loads, it returns the deformed shape and the internal force distribution.

\subsection{What Billow is not}

It contains \textbf{no aerodynamics}. There are no polars, no airfoils, no
lift model. Loads arrive as nodal forces from outside. Everything in this
document that involves pressure uses a prescribed uniform differential
pressure, not an aerodynamic solution. The same is true of \code{kite\_fem},
which is also purely structural.

It is also not a dynamics code: there are no inertial terms and no time
integration. Equilibrium is quasi-static.

\subsection{Isolation}

Billow depends only on NumPy and CasADi. It does not import the particle-system
solver (PSS), the vortex-step method (VSM), the aerostructural coupling loop or
any kite YAML schema, and nothing in those modules imports it. This is
deliberate and is the first invariant to preserve: it allows the element physics
to be validated against closed-form solutions before any of it touches a coupled
run. A coupling adapter, when written, belongs on the aerostructural side, not
inside Billow.

\subsection{Why a new model}

The existing minimum-energy solver in the AWETrim tree
(\code{aerostructural/pss/structural\_nlp.py}) already solves the current bridle
well. It does not extend, for two independent reasons.

\begin{description}[style=unboxed,leftmargin=1.2em]
\item[Physics.] Its element vocabulary is axial springs. Battens and
leading-edge tubes need bending and torsion; the canopy needs a membrane that
cannot carry compression.
\item[Scale.] It builds the objective with a Python loop over elements in
CasADi \code{MX}, so the expression graph grows with the element count. At tens
of springs this is free. At tens of thousands of fabric triangles, graph
construction --- and especially \code{gradient} and \code{hessian} over it ---
becomes the bottleneck long before the numerical solver does.
\end{description}

Billow keeps the formulation (still an NLP, still IPOPT, still minimum potential
energy) and replaces the assembly.

%======================================================================
\section{Formulation}

\subsection{Total potential energy}

The unknown configuration $X$ minimises
\begin{equation}
  \Pi(X) \;=\; \sum_{\text{sets}} U_{\text{set}}(X)
  \;-\; \mathbf{f}_{\text{ext}}\!\cdot\!\mathbf{x}
  \;-\; \mathbf{m}_{\text{ext}}\!\cdot\!\boldsymbol{\theta}(\boldsymbol{\psi})
  \;+\; \tfrac{1}{2}k_a\lVert \mathbf{x}-\mathbf{x}_{\text{seed}}\rVert^2 ,
  \label{eq:pi}
\end{equation}
subject only to bounds on pinned degrees of freedom. The last term is a weak
anchor discussed in \S\ref{sec:anchor}.

Because $\partial\Pi/\partial X = 0$ is the equilibrium equation, a converged
solve satisfies $\mathbf{f}_{\text{int}} + \mathbf{f}_{\text{ext}} = 0$ at every
free node to solver tolerance. There is no separate residual criterion to tune.

\subsection{Degrees of freedom}

\begin{equation}
  X = \big[\; \underbrace{\mathbf{x}_0 \ldots \mathbf{x}_{n-1}}_{\text{3 per node}}
  \;\big|\;
  \underbrace{\boldsymbol{\psi}_0 \ldots \boldsymbol{\psi}_{r-1}}_{\text{3 per rotational node}} \;\big]
\end{equation}

Three translations for \emph{every} node, and three incremental rotations only
for the $r$ nodes that some beam element needs a material frame for. Cable,
pulley and membrane nodes stay at three degrees of freedom. A canopy of 10\,000
fabric nodes therefore costs 30\,000 unknowns, not 60\,000, regardless of what
beams are attached elsewhere. For contrast, \code{kite\_fem} allocates six
degrees of freedom per node everywhere and masks the unused ones.

\subsection{Finite rotations: the Cayley parameterisation}

Rotations use the \textbf{Rodrigues (Cayley) vector}
$\boldsymbol{\psi} = 2\tan(\phi/2)\,\mathbf{n}$ rather than the exponential
map's rotation vector $\phi\,\mathbf{n}$. With $S = \operatorname{skew}(\boldsymbol{\psi})$,
\begin{align}
  R(\boldsymbol{\psi}) &= I + \frac{4}{4+\boldsymbol{\psi}\cdot\boldsymbol{\psi}}
    \left(S + \tfrac{1}{2}S^2\right), \label{eq:cayley}\\
  \boldsymbol{\psi}(R) &= \frac{2\,\operatorname{axial}(R-R^{\mathsf T})}{1+\operatorname{tr}R},
    \label{eq:cayleyinv}\\
  \boldsymbol{\psi}_{1/2} &= \frac{\boldsymbol{\psi}}
    {1+\sqrt{1+\tfrac{1}{4}\boldsymbol{\psi}\cdot\boldsymbol{\psi}}},
    \qquad R(\boldsymbol{\psi}_{1/2})^2 = R(\boldsymbol{\psi}).
    \label{eq:cayleyhalf}
\end{align}

The motivation is entirely numerical. Every map above is a \emph{rational}
function of its argument. There is no $\sin\phi/\phi$ removable singularity, so
no \code{if\_else} branch is needed at $\phi=0$, so the Hessian that IPOPT
differentiates is continuous everywhere in $|\phi| < \pi$. The half-angle
identity \eqref{eq:cayleyhalf} is worth noting: the naive form
$\tan(\phi/4) = (\sqrt{1+t^2}-1)/t$ has a removable $0/0$ at $t=0$;
rationalising it removes the singularity exactly.

\subsection{Incremental rotations and frame updates}

$\boldsymbol{\psi}_j$ is measured against a stored reference frame, so the
configuration is $R_j = R(\boldsymbol{\psi}_j)\,R_{\text{ref},j}$ and every
solve starts from $\boldsymbol{\psi} = 0$. After each solve the increment is
absorbed, $R_{\text{ref}} \leftarrow R(\boldsymbol{\psi}^\star) R_{\text{ref}}$,
and re-projected onto $SO(3)$ so repeated absorption cannot drift.

The Cayley chart is singular at $\phi = \pi$. Incremental rotations mean a node
can turn arbitrarily far \emph{across a load ramp} while each individual solve
stays well inside the chart. This is why applied moments require load stepping
(\S\ref{sec:moments}); nodal forces do not.

\subsection{Applied moments are not conjugate to \texorpdfstring{$\boldsymbol{\psi}$}{psi}}
\label{sec:moments}

A moment does work against the \emph{exponential-map} rotation vector, not
against the Rodrigues vector the degrees of freedom carry. The external work
term therefore converts,
\begin{equation}
  \boldsymbol{\theta}(\boldsymbol{\psi}) = \boldsymbol{\psi}\,
  \frac{2\arctan\!\left(\lVert\boldsymbol{\psi}\rVert/2\right)}{\lVert\boldsymbol{\psi}\rVert},
\end{equation}
which tends to $\boldsymbol{\psi}$ as $\lVert\boldsymbol{\psi}\rVert \to 0$.

This is not cosmetic. Using $-\mathbf{m}\cdot\boldsymbol{\psi}$ directly is
correct only to first order, and the roll-up benchmark
(\S\ref{sec:rollup}) exposed it as a \emph{load-path dependence}: the same
static problem gave $8.1\times10^{-3}L$ with four load steps and
$2.1\times10^{-1}L$ with eight. Statics must be path independent. With the
conversion in place the answer is identical to five digits across 4, 8 and 16
steps, and this is guarded by a regression test.

Two limitations remain, both fundamental rather than implementation artefacts:

\begin{itemize}
\item Applied moments need load stepping, because the chart is singular at
$\phi = \pi$. Four steps suffice for a full $2\pi$ roll-up.
\item A constant (``dead'') moment in three dimensions is \textbf{genuinely
non-conservative} and therefore has no potential at all. Energy minimisation can
represent only the fixed-axis case.
\end{itemize}

Nodal \emph{forces} --- which is what an aerodynamic coupling applies --- carry
none of these caveats and are exact in a single solve.

\subsection{Assembly: one mapped kernel per element type}
\label{sec:assembly}

This is the single design decision that makes the model scale. Each element
kernel is compiled once to a small CasADi \code{SX} function over \emph{one}
element and evaluated across the whole set with \code{Function.map} on a
gathered index matrix:
\[
  Q = \operatorname{reshape}\big(X[\,\mathrm{idx}\,],\; n_{\text{dof/el}},\; n_{\text{el}}\big),
  \qquad
  U_{\text{set}} = \textstyle\sum_e \mathcal{U}(Q_{:,e}, F_{:,e}, P_{:,e}).
\]
The objective graph therefore holds one node per element \emph{type}, not per
element. Graph construction, gradient and Hessian generation cost
$O(\text{number of element types})$, not $O(\text{number of elements})$.

The Hessian retains the sparsity of the element connectivity graph, so IPOPT's
sparse linear algebra sees a mesh-stencil matrix and scales accordingly. The
measured consequence is Table~\ref{tab:scaling}.

\begin{table}[H]
\centering
\caption{Flat clamped canopy, uniform out-of-plane load. Iteration count is
essentially mesh independent; cost sits in the sparse linear algebra.
Source: \code{run\_demo\_cases.py --case scaling}.}
\label{tab:scaling}
\begin{tabular}{rrrrr}
\toprule
Triangles & DOF & Build (s) & Solve (s) & IPOPT iterations \\
\midrule
128  & 243  & 0.07 & 0.05 & 6 \\
1800 & 2883 & 0.67 & 1.25 & 9 \\
5832 & 9075 & 3.0  & 6.1  & 9 \\
\bottomrule
\end{tabular}
\end{table}

\subsection{Everything numeric is a parameter}

Reference frames, every element parameter column (rest lengths, stiffnesses,
section properties), the external loads and the anchor are NLP
\emph{parameters}, not baked constants. One solver build therefore serves an
entire sweep: tape actuation, a stiffness ramp or a load sweep change parameter
values, never topology.

\subsection{The anchor term}
\label{sec:anchor}

The final term of \eqref{eq:pi} is a weak spring to the seed positions, present
only to pin a node that no taut element reaches and whose equilibrium position
would otherwise be indeterminate. It biases the reported force balance by at
most $k_a\lVert\mathbf{x}-\mathbf{x}_{\text{seed}}\rVert$; the default
$k_a = 10^{-6}$\,N/m puts that far below any coupling tolerance.

\subsection{Convergence criterion}

\code{converged} is the \emph{physical} verdict: the solve is accepted when the
free nodes are in force balance to
\[
  \max\big(\varepsilon_{\text{abs}},\; \varepsilon_{\text{rel}}\max_i\lVert \mathbf{f}_{\text{ext},i}\rVert\big),
  \qquad \varepsilon_{\text{abs}} = 10^{-6}\,\text{N},\;\; \varepsilon_{\text{rel}} = 10^{-6},
\]
whether or not IPOPT liked its own termination test. The relative term matters:
an absolute newton means very different things on a 1\,N bridle line and a
600\,N benchmark beam. IPOPT's own verdict is preserved alongside as
\code{ipopt\_success} and \code{status} --- on stiff, nearly quadratic
structures it routinely reports \code{Search\_Direction\_Becomes\_Too\_Small} at
an answer already exact to machine precision.

%======================================================================
\section{Element library}

Every kernel returns the strain energy of one element. Kernels are pure: no
loops over elements, no global indexing, no closing over per-element data.
Everything element-specific arrives through the parameter row.

\subsection{Cable}

A linear axial spring with stretch $s = \ell - \ell_0$:
\begin{equation}
  U = \tfrac{1}{2}k\,\langle s\rangle^2,
  \qquad \langle s\rangle = \max(0, s)\ \text{when tension-only}.
\end{equation}
The tension-only cut is $C^1$ but not $C^2$: the stiffness jumps across the
taut/slack boundary. An optional \code{slack\_smoothing} width rounds the corner
to recover a continuous Hessian at the cost of a small spurious compressive
stiffness.

\subsection{Pulley}

A rope running over a frictionless sheave at the middle node, with a single
shared stretch across both arms:
\begin{equation}
  U = \tfrac{1}{2}k\,\big\langle \lVert\mathbf{x}_2-\mathbf{x}_1\rVert
    + \lVert\mathbf{x}_3-\mathbf{x}_2\rVert - \ell_0 \big\rangle^2 .
\end{equation}
One shared stretch is exactly the statement that rope tension is equal either
side of the sheave, and it lets the pulley node find its own position along the
rope.

\paragraph{Convention warning.} \code{rest\_length} here is the \textbf{whole
rope}, both arms. This is not universal: the PSS geometry reader splits $\ell_0$
across the two arms, while \code{kite\_fem} stores the total on each arm. Any
geometry adapter must state which convention it is reading.

\subsection{Geometrically exact Timoshenko beam}

Two-node Simo--Reissner beam with one-point (mid-element) integration. Each node
carries a position and an orthonormal material frame $R = [\mathbf{d}_1\,
\mathbf{d}_2\, \mathbf{d}_3]$: twelve degrees of freedom per element.

With $\boldsymbol{\psi}$ the relative rotation between the two node frames from
\eqref{eq:cayleyinv}, and $R_m$ the exact half-way frame from
\eqref{eq:cayleyhalf}:
\begin{align}
  \boldsymbol{\Gamma} &= R_m^{\mathsf T}(\mathbf{x}_2-\mathbf{x}_1)/L_0,
    \label{eq:beamgamma}\\[2pt]
  \boldsymbol{\Omega} &= R_m^{\mathsf T}\boldsymbol{\psi}/L_0,
    \label{eq:beamomega}\\[2pt]
  U &= \frac{L_0}{2}\Big[
     \Delta\boldsymbol{\Gamma}^{\mathsf T}\operatorname{diag}(EA, GA_2, GA_3)\,\Delta\boldsymbol{\Gamma}
   + \Delta\boldsymbol{\Omega}^{\mathsf T}\operatorname{diag}(GJ, EI_2, EI_3)\,\Delta\boldsymbol{\Omega}
    \Big].
    \label{eq:beamenergy}
\end{align}

$\boldsymbol{\Gamma}$ is a \emph{deformation} measure rather than a strain: its
first component is the axial stretch, which is unity in the reference
configuration, and the other two are the transverse shears.
$\boldsymbol{\Omega}$ holds the torsion and the two bending curvatures. What
drives the energy is the change against the reference,
$\Delta(\cdot) = (\cdot) - (\cdot)_0$, so an initially curved or pre-twisted
member --- a leading-edge tube, a strut --- is stress-free as built.

Three properties make this element usable inside an NLP objective:

\begin{description}[style=unboxed,leftmargin=1.2em]
\item[Objective.] Under a superposed rigid rotation $Q$, both $\boldsymbol{\psi}$
and $\mathbf{x}_2-\mathbf{x}_1$ map to $Q(\cdot)$ while $R_m \to QR_m$, so
$\boldsymbol{\Gamma}$ and $\boldsymbol{\Omega}$ are invariant and rigid-body
motion costs exactly zero energy. A slack member that swings freely is not
spuriously stiffened.
\item[Shear-locking free.] The single midpoint sample is the standard reduced
integration cure. Kite battens and leading-edge tubes are slender, so this is
not optional; \S\ref{sec:locking} verifies it at slenderness 800.
\item[Branch-free.] Using the Cayley vector as the curvature measure makes the
element strain $\boldsymbol{\Omega}_{\text{exact}} + O(\phi^3)$ per element ---
the same order as the one-point quadrature error, and vanishing under mesh
refinement. This is a deliberate trade: the exponential map's $\sin\phi/\phi$
would need a branch at $\phi=0$.
\end{description}

\subsection{Inflatable tube}
\label{sec:inflatable}

An inflated leading-edge tube is not linear-elastic. Its bending stiffness falls
away as the compressed side wrinkles, until the section carries a limiting
moment and then collapses. \code{kite\_fem} captures this with empirical fits
(coefficients $C_1 \ldots C_{19}$, from the ASKITE work) applied as \emph{secant}
stiffnesses: $EI$ and $GJ$ are recomputed each iteration from the current
deflection and twist.

That is consistent inside a Newton residual, but it is \textbf{not a potential}.
Substituting a state-dependent $EI(\kappa)$ into $\tfrac{1}{2}EI\kappa^2$
silently drops the $\mathrm{d}EI/\mathrm{d}\kappa$ terms, so the gradient would
no longer be the internal moment and a minimum-energy solve would converge to
the wrong equilibrium. The fits must be \emph{integrated}, not substituted.

\paragraph{From tip-load fit to moment--curvature law.}
The bending fit is published as tip load against normalised tip deflection
$v = \delta/L$ of a \textbf{one-metre} cantilever:
$P(v) = D\big(1-e^{-(N/D)v}\big)$. For a linear cantilever
$\delta = PL^3/3EI$ and $\kappa_{\text{root}} = PL/EI$, so $\kappa = 3v/L$, and
the root moment is $M = PL$. At the calibration length $L = 1$\,m this converts
the fit into an intrinsic constitutive law:
\begin{align}
  M(\kappa) &= M_{\max}\left(1 - e^{-EI_0\kappa/M_{\max}}\right),
   \qquad EI_0 = N/3, \quad M_{\max} = D, \label{eq:tubemoment}\\
  W_b(\kappa) &= M_{\max}\left[\,|\kappa| - k_0\left(1 - e^{-|\kappa|/k_0}\right)\right],
   \qquad k_0 = M_{\max}/EI_0 . \label{eq:tubeenergy}
\end{align}
$W_b$ is quadratic ($\tfrac{1}{2}EI_0\kappa^2$) near zero, so it is smooth there
despite the $|\kappa|$, and saturates at $M_{\max}$.

Recasting this way also removes a length inconsistency: \code{kite\_fem} infers
$EI = P/(3v)$, which equals the true $EI$ only for a one-metre element, because
$v$ is already normalised by element length. A moment--curvature law is
length-independent by construction.

\paragraph{Torsion} is already intrinsic (twist per unit length) and integrates
directly:
\begin{equation}
  T(\omega) = c_1\arctan(c_2\omega),
  \qquad
  W_t(\omega) = c_1\!\left[\omega\arctan(c_2\omega)
    - \frac{\ln\!\left(1+c_2^2\omega^2\right)}{2c_2}\right],
\end{equation}
with $GJ_0 = c_1c_2$ and a limiting torque $c_1\pi/2$.

Biaxial bending uses the resultant curvature $\sqrt{\Omega_2^2+\Omega_3^2}$, the
correct isotropic extension for an axisymmetric tube. Axial and shear stiffness
are \emph{not} covered by the fits and stay linear; supply them from the tube
geometry.

\begin{table}[H]
\centering
\caption{Derived tube properties from the ASKITE fits.
Higher pressure stiffens and strengthens; the thinner tube collapses
\emph{before} its moment saturates ($\kappa_{\text{collapse}} < k_0$).}
\label{tab:tube}
\begin{tabular}{llrrrrrr}
\toprule
$d$ (m) & $p$ (bar) & $EI_0$ (N\,m$^2$) & $M_{\max}$ (N\,m) & $k_0$ (m$^{-1}$)
 & $\kappa_{\text{collapse}}$ (m$^{-1}$) & $GJ_0$ (N\,m$^2$) & $T_{\max}$ (N\,m)\\
\midrule
0.16 & 0.3 & 288.2 & 58.10 & 0.2016 & 0.2751 & 91.3 & 124.9 \\
0.16 & 0.5 & 362.0 & 98.78 & 0.2728 & 0.2973 & 86.5 & 174.7 \\
0.12 & 0.3 & 126.8 & 34.16 & 0.2694 & 0.2217 & 33.2 & 117.1 \\
\bottomrule
\end{tabular}
\end{table}

\paragraph{Collapse is reported, not enforced.}
A genuinely dropping post-collapse moment would make the energy \emph{decrease}
with curvature --- an unbounded mechanism that minimisation would simply run
away from. The law therefore stops at saturation, and
\code{inflatable\_beam\_state} returns a \code{utilisation} ratio
$\kappa/\kappa_{\text{collapse}}$ and a \code{collapsed} flag, so leaving the
calibrated range is visible rather than silent.

\subsection{Wrinkling membrane}
\label{sec:membrane}

Constant-strain triangle with a Saint-Venant--Kirchhoff plane-stress law, in the
finite-strain (Green--Lagrange) measure so a canopy panel can rotate and billow
arbitrarily without spurious stiffening. Translation-only: a membrane carries no
bending.

The reference triangle is flattened once into its own in-plane basis and stored
as the inverse edge matrix $D_0^{-1}$, so the kernel only ever sees a $2\times2$
map:
\begin{equation}
  F = [\,\mathbf{a}\;\mathbf{b}\,]\,D_0^{-1}, \quad
  C = F^{\mathsf T}F, \quad
  E = \tfrac{1}{2}(C-I),
\end{equation}
\begin{equation}
  W_{\text{taut}} = \frac{E_{\text{mod}}}{2(1-\nu^2)}
   \left[(\operatorname{tr}E)^2 - 2(1-\nu)\left(E_{11}E_{22}-E_{12}^2\right)\right].
\end{equation}

\paragraph{Wrinkling.} A membrane cannot carry compression. Left alone, the SVK
law lets the canopy mesh buckle into whatever compressive mode the
discretisation offers, and the result depends on mesh density rather than on
physics. The relaxed (tension-field) energy in principal Green strains
$e_1 \ge e_2$ is the classical Pipkin form:
\begin{equation}
  W_{\text{relaxed}} =
  \begin{cases}
    W_{\text{taut}}, & e_2 + \nu e_1 \ge 0 \quad\text{(taut)},\\[2pt]
    \tfrac{1}{2}E_{\text{mod}}\,e_1^2, & e_1 > 0 \ \text{otherwise}\quad\text{(wrinkled)},\\[2pt]
    0, & e_1 \le 0 \quad\text{(slack)}.
  \end{cases}
\end{equation}

This is exactly what a minimum-energy formulation wants: the relaxed functional
is the quasiconvex envelope, so the minimiser lands on the wrinkled state
directly instead of chasing the near-singular tangent a residual-form Newton
solve has to fight through. \textbf{Solving fabric by energy minimisation is not
a workaround --- it is the better-posed formulation.}

\paragraph{Slack fabric needs a residual stiffness.} A fully slack region stores
\emph{exactly} zero energy, so its Hessian block is exactly zero and the Newton
step is undefined; IPOPT reports \code{Error\_In\_Step\_Computation} on a panel
clamped into a frame smaller than itself. Blending a small fraction of the
unrelaxed law back in,
\begin{equation}
  W = (1-r)\,W_{\text{relaxed}} + r\,W_{\text{taut}}, \qquad r = 10^{-4},
\end{equation}
leaves the taut region untouched (there the two coincide) while giving slack
fabric a definite tangent. It is also not a fiction: real fabric has some
compressive and bending resistance. Do not set $r=0$ on a canopy that can go
slack.

%======================================================================
\section{Validation}
\label{sec:validation}

All references in this section come from \emph{outside} the code: an exact
boundary-value solution, an exact circle, published tabulated values, or the
independent fits themselves. Nothing here can drift with the implementation.

Where \code{kite\_fem} appears it is driven with matched linear section
properties and its inflatable constitutive update disabled, so the two share one
material model. Section properties are deliberately slender: the classical
references are inextensible and shear-rigid, so axial and shear flexibility are
kept below $\sim0.1\%$ and cannot be mistaken for discretisation error.

\subsection{Euler elastica}

The \emph{elastica} is the classical large-deflection equilibrium shape of a
thin, inextensible, shear-rigid rod --- the problem Jacob Bernoulli posed and
Euler solved in 1744, whose exact solution is expressible in elliptic integrals.
It is a \textbf{reference solution, not a beam formulation}: it generalises
nothing about the element used here, and no beam element is involved in
producing it. Its value as a benchmark is that it stays exact at arbitrarily
large rotation, which is precisely where a geometrically exact element has to be
tested and where a linearised one fails.

The case is a tip-loaded cantilever at load parameter $\alpha = PL^2/EI$ up to
10, which takes the rod through most of a quarter turn. The reference comes from
solving the elastica equation $\theta'' = -\alpha\cos\theta$, $\theta(0)=0$,
$\theta'(1)=0$ as a boundary-value problem to $10^{-10}$ with continuation in
$\alpha$.

\begin{figure}[H]
\centering
\includegraphics[width=\textwidth]{elastica.png}
\caption{Tip locus and error against the exact elastica, 40 elements.
Billow tracks the exact curve across the whole range; \code{kite\_fem}
departs progressively and its tip deflection saturates beyond $\alpha \approx 2$.}
\label{fig:elastica}
\end{figure}

\begin{table}[H]
\centering
\caption{Tip position error normalised by length, 40 elements.}
\begin{tabular}{lrrrr}
\toprule
$\alpha$ & 0.5 & 2 & 6 & 10 \\
\midrule
Billow           & $9.0\times10^{-5}$ & $2.4\times10^{-4}$ & $6.5\times10^{-4}$ & $1.1\times10^{-3}$\\
\code{kite\_fem} & $1.6\times10^{-2}$ & $1.8\times10^{-1}$ & $4.2\times10^{-1}$ & $4.9\times10^{-1}$\\
\bottomrule
\end{tabular}
\end{table}

\paragraph{Mesh refinement separates discretisation from formulation.}
At $\alpha = 3$, refining from 10 to 160 elements:

\begin{table}[H]
\centering
\begin{tabular}{lrrrrr}
\toprule
Elements & 10 & 20 & 40 & 80 & 160\\
\midrule
Billow           & $3.8\times10^{-3}$ & $7.0\times10^{-4}$ & $3.1\times10^{-4}$
                 & $4.6\times10^{-4}$ & $5.0\times10^{-4}$\\
\code{kite\_fem} & $2.908\times10^{-1}$ & $2.914\times10^{-1}$ & $2.916\times10^{-1}$
                 & $2.917\times10^{-1}$ & $2.917\times10^{-1}$\\
\bottomrule
\end{tabular}
\end{table}

The \code{kite\_fem} error is \textbf{mesh independent}, which places it in the
formulation rather than the discretisation. Its beam matches linear theory to
$0.02\%$ below about one degree of tip rotation, so the property matching used
here is correct; the error is purely rotation-driven. Envelope: below $1\%$ up
to roughly $10^\circ$, about $15\%$ at $\alpha=1$, saturating beyond $\alpha=2$.
It reports \code{converged=True} at those answers.

\subsection{Roll-up under an end moment}
\label{sec:rollup}

A cantilever under $M = 2\pi EI/L$ must roll into an exact closed circle, so the
reference requires no lookup at all.

\begin{figure}[H]
\centering
\includegraphics[width=\textwidth]{rollup.png}
\caption{Roll-up at 30, 60 and 100\% of $2\pi EI/L$, and the distance from the
exact circular arc. Billow converges at second order; \code{kite\_fem} moves
away from the exact answer under refinement.}
\label{fig:rollup}
\end{figure}

\begin{table}[H]
\centering
\caption{Tip distance from the exact arc, normalised by length.}
\begin{tabular}{lrrr}
\toprule
Elements & 10 & 20 & 40 \\
\midrule
Billow (4 load steps) & $1.13\times10^{-1}$ & $3.16\times10^{-2}$ & $8.14\times10^{-3}$\\
\code{kite\_fem} (single shot) & $5.90\times10^{-2}$ & $1.71\times10^{-1}$ & $2.32\times10^{-1}$\\
\bottomrule
\end{tabular}
\end{table}

Billow's error ratios are $3.6$ and $3.9$, i.e.\ second order. Extending to 80
elements gives $2.05\times10^{-3}$, confirming the trend. This benchmark is what
exposed the moment-conjugacy defect of \S\ref{sec:moments}.

\subsection{Bathe--Bolourchi \texorpdfstring{$45^\circ$}{45-degree} bend}

The standard three-dimensional finite-rotation benchmark: a $45^\circ$ circular
arc of radius 100, square section, out-of-plane tip load $F = 600$. Published
tip \emph{displacements} differ by about $1\%$ between authors, so the reference
is a band.

\begin{figure}[H]
\centering
\includegraphics[width=\textwidth]{bend45.png}
\caption{Tip displacement against published references, and convergence.}
\label{fig:bend45}
\end{figure}

\begin{table}[H]
\centering
\caption{Tip displacement, 32 elements. The arc here starts at the origin with
its tangent along $+x$ and sweeps towards $+y$, which swaps the two in-plane
components relative to the usual tabulation.}
\begin{tabular}{lrrrr}
\toprule
& $u_x$ & $u_y$ & $u_z$ & $\lVert\mathbf{u}\rVert$\\
\midrule
Simo \& Vu-Quoc (1986) & $-23.48$ & $-13.50$ & $53.37$ & $59.85$\\
Bathe \& Bolourchi (1979) & $-23.5$ & $-13.4$ & $53.4$ & $59.85$\\
Crisfield (1990) & $-23.5$ & $-13.6$ & $53.4$ & $59.88$\\
\midrule
\textbf{Billow} & $\mathbf{-23.52}$ & $\mathbf{-13.60}$ & $\mathbf{53.44}$ & $\mathbf{59.95}$\\
\code{kite\_fem} & $-37.74$ & $-22.14$ & $28.17$ & $52.04$\\
\bottomrule
\end{tabular}
\end{table}

Billow sits within the spread between published references ($0.17\%$ on the
magnitude). \code{kite\_fem} is $47\%$ low out of plane.

\subsection{Inflatable tube law}

The port of \S\ref{sec:inflatable} is verified with \textbf{pure end moments},
where the exact solution is a constant-curvature arc at whatever curvature the
fit prescribes. That check is exact at any deflection --- unlike the tip-load
form the fit was published in, whose own inversion assumes linear cantilever
theory and therefore cannot serve as a verification instrument.

\begin{figure}[H]
\centering
\includegraphics[width=\textwidth]{inflatable.png}
\caption{ASKITE fits (lines) with Billow end-moment solves on top (markers).
Hollow markers lie beyond the collapse curvature, where the fit is
extrapolation: the solve is fine, the calibration is not.
This figure completes \code{kite\_fem/examples/FEM\_beam\_verification.py},
which plots the fitted curves but leaves the simulation loop as comments.}
\label{fig:inflatable}
\end{figure}

Worst errors: $0.031\%$ (bending, 0.3\,bar), $0.057\%$ (bending, 0.5\,bar),
$0.749\%$ and $1.491\%$ (torsion).

\subsection{Beam convergence and shear locking}
\label{sec:locking}

Tip deflection of a slender cantilever must approach
$PL^3/3EI + PL/\kappa GA$ at second order:

\begin{table}[H]
\centering
\begin{tabular}{lrrrrrr}
\toprule
Elements & 2 & 4 & 10 & 20 & 40 & 80\\
\midrule
Error vs Timoshenko & $-6.25\%$ & $-1.56\%$ & $-0.250\%$ & $-0.063\%$
 & $-0.016\%$ & $-0.004\%$\\
\bottomrule
\end{tabular}
\end{table}

The error falls by a factor of four per refinement. Locking is checked
separately by holding the load and lengthening the beam through slenderness
ratios of 50, 200 and 800: the tip deflection tracks beam theory to $0.2\%$
throughout, where a locking element would stiffen towards zero.

\subsection{Membrane}

Checked against closed form rather than against another code:

\begin{itemize}
\item \textbf{Rigid-body motion} costs zero energy under arbitrary rotation and
translation --- what the finite-strain measure buys over a small-strain one.
\item \textbf{Equibiaxial stretch} matches the analytic SVK plane-stress energy
to $10^{-9}$ relative.
\item \textbf{Wrinkled branch} is exactly uniaxial: stretching along $x$ while
squashing along $y$ gives $\tfrac{1}{2}E t A e_1^2$, and squashing further
changes nothing.
\item \textbf{Slack panel} stores exactly zero, where the unrelaxed law would
store large energy.
\item \textbf{Flat-plate inflation} converges under mesh refinement from a
configuration with zero out-of-plane stiffness --- the singular start that a
residual-form Newton method has nothing to invert at.
\end{itemize}

\subsection{Test suite}

86 tests, all passing, organised one file per concern:

\begin{table}[H]
\centering
\begin{tabular}{ll}
\toprule
\code{test\_rotations.py} & Cayley round trips, half-angle identity\\
\code{test\_cable.py} & Hooke parity with PSS, slack cut, pulley tension equalisation\\
\code{test\_beam.py} & objectivity, $O(h^2)$ convergence, no locking, torsion vs $GJ$\\
\code{test\_membrane.py} & analytic SVK, three wrinkling branches, inflation, scale\\
\code{test\_inflatable.py} & energy integrates the fitted moment; end-moment solves\\
\code{test\_benchmarks.py} & roll-up vs exact circle, load-step independence, bend45\\
\code{test\_energy.py} & mapped assembly vs a naive per-element sum; DOF layout\\
\bottomrule
\end{tabular}
\end{table}

The most important of these is in \code{test\_energy.py}: the total energy
computed through the gather-and-map machinery is compared against the same
energy assembled one element at a time in plain Python, on a deliberately
irregular model containing every element type with a beam running over a
scattered, out-of-order subset of the node numbering. That machinery is easy to
get subtly wrong --- a column-major reshape, a rotational slot off by one ---
and wrong in a way that still converges to a plausible-looking shape.

%======================================================================
\section{Validation against a measured kite}
\label{sec:hanging}

The benchmarks of Section~\ref{sec:validation} are exact solutions of idealised
problems. This section is the only place where Billow meets a real structure
with measured deformation: the TU~Delft V3 suspended from a bridle bar in a
hangar, at two inflation pressures and five loadings, with its deformed shape
recorded by stereoscopic photogrammetry.

Measurements, the reference model and \code{kite\_fem}'s published results are
vendored under \code{data/LEI-V3-KITE/hanging\_validation/} from
\code{awegroup/kite\_fem} (MIT, \textcopyright~2025 Patrick Roeleveld),
\code{main} at \code{630c8c0}. The three-dimensional marker clouds come from
Haanen's repository.\footnote{%
\url{https://github.com/pimjhaanen/photogrammetry_thesis}}
Thesis references are to P.\,I.\,H.~Roeleveld, \emph{Fast finite element
modelling of bridled leading-edge inflatable kites}, MSc thesis, TU~Delft,
February 2026. Everything here is regenerated by the \code{run\_hanging\_*}
scripts in \code{scripts/structural/}.

Three results carry beyond this experiment. The distributed reference model is
\emph{ill posed} --- its anchor node belongs to no element --- so any future use
of the dataset must supply a suspension. \code{kite\_fem}'s beams converge to
about $0.23\times$ its own constitutive law, so it is not a clean reference
solution. And the dataset is governed almost entirely by inflatable-tube
bending, not by the canopy it was chosen to test.

\subsection{The experiment}

The kite hangs from its bridle and the deformed shape is recorded
stereoscopically, with markers on the inflatable structure. Twelve lengths are
extracted per load case: nine trailing-edge segments \texttt{La}--\texttt{Li}
between consecutive struts, the tip-to-tip span \texttt{b}, and two
tip-to-centre-leading-edge diagonals \texttt{LcsTL}, \texttt{LcsTR}.

The deformation is large. Under self-weight alone the span falls from
\SI{8.305}{m} as built to a measured \SI{4.867}{m} --- a 41\% closure --- and
every trailing-edge segment shortens by 26--43\%. The tip-load cases run the
same deformation backwards: spreading the tips with \SI{2}{kg} or \SI{5}{kg}
reopens the span to \SI{7.8}{}--\SI{8.3}{m}, close to the as-built shape.

\paragraph{Suspension.} The kite hangs from a \textbf{bar}, not a point.
Roeleveld's thesis (Appendix~B, Table~B.1) marks \texttt{A\_main},
\texttt{A\_I}, \texttt{M}, \texttt{PL} and \texttt{SL} as \emph{N/A} for the
hanging test and adds a \texttt{bar} of \SI{1650}{mm}; Section~6.3 explains
that the bridle was adapted to fit the hangar height. The same table records
the one documented shortening, \texttt{B\_r,5} from \SI{12168}{}
to~\SI{9425}{mm}.

\paragraph{Provenance caveat.} The upstream repository describes the
measurements only as ``experimental validation reference data'' and never uses
the word photogrammetry; the thesis identifies the method as stereoscopic
camera measurement. The published rows are also perfectly symmetric
left-to-right to four decimals, so they have been mirrored or
port/starboard-averaged.

\subsection{Rebuilding the model}

Billow is rebuilt from \texttt{hanging\_test\_initial.npz}, the fully assembled
reference model, rather than from the YAML --- there is no committed
YAML$\rightarrow$npz builder, so the \texttt{.npz} \emph{is} the input both
codes solve. Nodes, connectivity, rest lengths, stiffnesses and tube diameters
therefore transfer with no re-interpretation:

\begin{center}
\begin{tabular}{ll}
\toprule
shared, element for element & 276 nodes, 46 bridle cables, 14 pulleys,
                              98 inflatable tube beams \\
changed & canopy: 721 tension-only springs $\rightarrow$ 392 wrinkling
          membrane triangles \\
\bottomrule
\end{tabular}
\end{center}

Two checks make this trustworthy rather than assumed. The load vector
reproduces \texttt{kite\_fem}'s own \texttt{loading()} to \textbf{exactly zero}
for all five load combinations, and the length extractor reproduces its
published \texttt{model\_results.csv} from its saved states to
\textbf{machine zero}.

\paragraph{Canopy grid.} The reference canopy is a structured $29\times8$ grid
stored only as a flat spring list. It is recovered geometrically --- the
reference builder projects every interior section node onto the straight line
joining its leading- and trailing-edge node, so each chordwise section is
exactly collinear --- and then \emph{verified}: all 29 sections must come out
with 8 nodes, and all 721 canopy springs must be grid edges with no spring
joining two grid nodes being anything else. Both hold exactly.

\subsection{Defects found in the reference model}

\subsubsection{The anchor is connected to nothing}

Node~0 is flagged fixed but appears in no spring, no pulley and no beam; seven
further bridle nodes are orphaned with it (all of exactly zero mass). The kite
is therefore a free body under \SI{126}{N} of unbalanced weight, and its
potential is unbounded below. Billow reports that faithfully: the structure
translates hundreds of kilometres, and neither load stepping nor constitutive
continuation helps, because none of them supplies the missing constraint.

\texttt{kite\_fem} does not diverge only because its \texttt{I\_stiffness}
regularisation acts as a ground spring on every node. Its own published
solutions still drift \SIrange{0.29}{0.34}{m} bodily upward, with node~0 sitting
fixed and unconnected throughout.

\subsubsection{Bridle lengths disagree with the thesis}

Twenty-eight of the thirty-six declared bridle lines match Table~B.1 to the
millimetre, which is what makes the remaining eight look like transcription
errors rather than naming mismatches: \texttt{BRI} and \texttt{BRII} are an
exact swap; \texttt{abcd2} and \texttt{abcd3} carry a duplicated leading digit;
\texttt{ab3} equals \texttt{ab2}; \texttt{a5} and \texttt{br5} differ; and
\texttt{brmain ext} is present in the table but unwired.

The eighth is not a transcription error but a contradiction inside the thesis
itself. Table~B.1 marks \texttt{A\_I} as N/A for the hanging test, while
Figure~B.1 \emph{draws} \texttt{A\_I} running to the bar in the adapted
geometry. The two readings are structurally different --- cutting the line
removes a front load path --- so both were run:

\begin{table}[htbp]
\centering
\caption{Three readings of the same source. Cutting \texttt{A\_I} is the most
accurate reading on the cases the model can solve and the worst overall,
because removing that load path is what lets the four centre-load cases fall
into a folded local minimum. No reading is adopted here: the ambiguity is a
finding about the source data, not something to settle by taking whichever
scores best.}
\begin{tabular}{lrrrr}
\toprule
bridle reading & all ten & six solvable & four centre-load & folded \\
 & (\%) & (\%) & (\%) & cases \\
\midrule
bridle as stored & 15.1 & 11.0 & 21.3 & 0 \\
Table~B.1, $A_I$ cut & 17.3 & 8.9 & 30.0 & 4 \\
Table~B.1, $A_I$ kept & 15.7 & 11.5 & 21.9 & 0 \\
\texttt{kite\_fem} & 12.4 & 7.8 & 19.3 & 0 \\
\bottomrule
\end{tabular}
\end{table}

\subsubsection{The reference beams converge far softer than their own law}

Billow's integrated tube law agrees with \texttt{kite\_fem}'s secant $EI$ to
1--3\% across the whole curvature range, departing only below its
\num{0.03} normalised-deflection floor. The two codes therefore implement the
same constitutive law. They do not converge to the same stiffness:

\begin{center}
\begin{tabular}{rrrrr}
\toprule
load case & median $EI$ (\si{N.m^2}) & vs.\ own law & min $EI$ & beams flagged collapsed \\
\midrule
1 & 24.5 & $0.29\times$ & 2.7 & 78 \\
4 & 69.4 & $0.81\times$ & $-42.4$ & 49 \\
6 & 43.1 & $0.50\times$ & $-164$ & 76 \\
9 & 84.1 & $0.98\times$ & $-459$ & 40 \\
\bottomrule
\end{tabular}
\end{center}

Three of the four cases contain \emph{negative} bending stiffness, which the
secant form $EI = PL^3/\bigl(3(v - PL/kGA)\bigr)$ produces as soon as the
deflection estimate falls below the shear term. \texttt{kite\_fem} therefore
lands closer to the measurement than Billow because its beams are several times
softer than the law it implements, not because its physics is closer.

\subsection{Results}

\begin{table}[htbp]
\centering
\caption{All ten load cases. Errors are the mean over the twelve measured
lengths of $|L_\text{model}-L_\text{meas}|/L_\text{meas}$; the reference
\texttt{shapecorrelation} metric is not used, because as a dot-product ratio
dominated by the two quantities of order \SIrange{4}{8}{m} it reads
0.980--0.9997 even where the span is 48\% wrong. $\|r\|$ is the force residual
on the worst free node.}
\label{tab:hang-results}
\begin{tabular}{rccc rrr rr rr}
\toprule
 & \multicolumn{3}{c}{load case} & \multicolumn{3}{c}{span (m)} & \multicolumn{2}{c}{mean rel.\ error (\%)} & \multicolumn{2}{c}{Billow} \\
\cmidrule(lr){2-4}\cmidrule(lr){5-7}\cmidrule(lr){8-9}\cmidrule(lr){10-11}
LC & $p$ (bar) & tip (kg) & centre (kg) & meas. & Billow & \texttt{kite\_fem} & Billow & \texttt{kite\_fem} & $t$ (s) & $\|r\|$ (N) \\
\midrule
1 & 0.15 & 0 & 0 & 4.867 & 7.276 & 5.382 & 25.9 & 13.7 & 7 & 2.2e-07 \\
2$^{\dagger}$ & 0.15 & 0 & 9.7 & 5.042 & 7.554 & 7.462 & 27.6 & 25.0 & 15 & 9.1e-08 \\
3$^{\dagger}$ & 0.15 & 0 & 25.2 & 5.113 & 7.841 & 7.576 & 21.8 & 18.9 & 19 & 3.8e-07 \\
4 & 0.15 & 2 & 0 & 7.843 & 8.597 & 8.467 & 6.5 & 6.7 & 16 & 1.1e-07 \\
5 & 0.15 & 5 & 0 & 8.271 & 8.916 & 8.962 & 5.9 & 5.4 & 18 & 4.1e-07 \\
6 & 0.25 & 0 & 0 & 5.453 & 7.468 & 6.363 & 19.5 & 9.3 & 6 & 4.0e-09 \\
7$^{\dagger}$ & 0.25 & 0 & 9.7 & 5.575 & 7.690 & 7.691 & 19.7 & 18.9 & 15 & 3.7e-09 \\
8$^{\dagger}$ & 0.25 & 0 & 25.2 & 5.697 & 7.965 & 7.638 & 18.4 & 14.6 & 17 & 6.4e-09 \\
9 & 0.25 & 2 & 0 & 7.818 & 8.593 & 8.461 & 5.2 & 5.6 & 15 & 1.1e-07 \\
10 & 0.25 & 5 & 0 & 8.224 & 8.908 & 8.955 & 6.2 & 5.9 & 16 & 4.0e-07 \\
\midrule
\multicolumn{7}{r}{mean, all ten} & 15.7 & 12.4 & 14 & \\
\multicolumn{7}{r}{mean, excluding centre-load cases} & 11.5 & 7.8 & 13 & \\
\bottomrule
\end{tabular}
\end{table}

\begin{figure}[htbp]
\centering
\includegraphics[width=\textwidth]{summary.png}
\caption{Span and mean relative error across the ten cases. Both codes
overstate the span wherever they solve the physical branch; the gap is widest
where gravity alone sets the shape and narrowest where an imposed tip load
dominates. The shaded centre-load cases are reproduced by neither code, and
are where Billow's folded branch shows as a near-zero span.}
\end{figure}

\begin{table}[htbp]
\centering
\caption{Per-quantity error, averaged over the ten cases. The two codes fail
differently, which the signed-bias columns separate from the magnitude
columns: Billow's bias is one-signed --- it overstates essentially every
length, i.e.\ its wing stays too open --- whereas \texttt{kite\_fem} scatters,
understating the outboard segments and the tip-to-centre diagonals while
overstating the inboard ones. A single stiffness factor could plausibly
account for the first and could not account for the second.}
\begin{tabular}{lrrrrr}
\toprule
quantity & meas.\ range (m) & \multicolumn{2}{c}{$|$rel.\ error$|$ (\%)} & \multicolumn{2}{c}{signed bias (\%)} \\
\cmidrule(lr){3-4}\cmidrule(lr){5-6}
 & & Billow & \texttt{kite\_fem} & Billow & \texttt{kite\_fem} \\
\midrule
\texttt{La} & 0.848--0.903 & 2.4 & 2.4 & +1.9 & -0.6 \\
\texttt{Lb} & 0.979--1.280 & 16.5 & 10.6 & +16.5 & +8.3 \\
\texttt{Lc} & 0.903--1.276 & 18.8 & 14.3 & +18.8 & +9.7 \\
\texttt{Ld} & 0.749--1.259 & 24.5 & 20.8 & +24.5 & +20.8 \\
\texttt{Le} & 0.920--1.217 & 13.0 & 16.2 & +13.0 & +16.2 \\
\texttt{Lf} & 0.749--1.259 & 24.1 & 20.8 & +24.1 & +20.8 \\
\texttt{Lg} & 0.903--1.276 & 18.4 & 14.3 & +18.4 & +9.7 \\
\texttt{Lh} & 0.979--1.280 & 16.9 & 10.6 & +16.9 & +8.3 \\
\texttt{Li} & 0.848--0.903 & 2.4 & 2.4 & +1.9 & -0.6 \\
\texttt{b} & 4.867--8.271 & 30.3 & 22.9 & +30.3 & +22.9 \\
\texttt{LcsTL} & 4.029--4.888 & 10.3 & 6.8 & +10.3 & +5.1 \\
\texttt{LcsTR} & 4.029--4.888 & 10.4 & 6.8 & +10.4 & +5.1 \\
\bottomrule
\end{tabular}
\end{table}

\begin{figure}[htbp]
\centering
\includegraphics[width=\textwidth]{lengths.png}
\caption{The twelve measured quantities per load case, model over measurement.}
\end{figure}

\subsubsection{The solved kite}

\begin{figure}[htbp]
\centering
\includegraphics[width=0.94\textwidth]{canopy_1.png}\\[-1mm]
\includegraphics[width=0.94\textwidth]{canopy_5.png}
\caption{The solved shape under self-weight (top) and with the tips spread by
\SI{5}{kg} (bottom), with the measured markers overlaid. The canopy is shaded
by facet orientation, so the billow between struts and the sag of each bay are
visible directly. Both are the $EI\times0.36$ solution.}
\end{figure}

\begin{figure}[htbp]
\centering
\includegraphics[width=0.94\textwidth]{regime_1.png}
\caption{The same self-weight case with each canopy triangle coloured by its
tension-field state. This is a prediction the membrane formulation makes and a
tension-only spring net cannot: under self-weight almost the whole canopy is
wrinkled --- one principal stress relaxed to zero --- with taut material only
where the bridle pulls through, and a handful of fully slack triangles at the
tip.}
\end{figure}

\clearpage
\subsubsection{Geometry, case by case}

The figures that follow show the inflatable structure --- leading edge and the
ten struts --- in front, top and side projection, with the as-built shape for
reference. The canopy is omitted for legibility; it carries little of the load
in this case (Section~\ref{sec:hang-scope}).

Black dots are Haanen's photogrammetry markers, carried into the model frame by
a rigid transform. For these figures the registration uses the \emph{leading
edge alone}, so the two leading edges lie on top of one another and the
remaining disagreement appears where it physically is --- in the struts and the
canopy --- rather than being spread evenly over the structure by a global fit.
The error table in Section~\ref{sec:hang-shape} uses its own per-model alignment over
all curves, which is the fair one for scoring; these overlays are for reading,
not for measuring. The markers stop short of the tips because the outermost bays
are not instrumented.

\begin{figure}[htbp]
\centering
\includegraphics[width=\textwidth]{geometry_1.png}
\caption{Load case 1.}
\end{figure}

\begin{figure}[htbp]
\centering
\includegraphics[width=\textwidth]{geometry_2.png}
\caption{Load case 2.}
\end{figure}

\begin{figure}[htbp]
\centering
\includegraphics[width=\textwidth]{geometry_3.png}
\caption{Load case 3.}
\end{figure}
\clearpage

\begin{figure}[htbp]
\centering
\includegraphics[width=\textwidth]{geometry_4.png}
\caption{Load case 4.}
\end{figure}

\begin{figure}[htbp]
\centering
\includegraphics[width=\textwidth]{geometry_5.png}
\caption{Load case 5.}
\end{figure}

\begin{figure}[htbp]
\centering
\includegraphics[width=\textwidth]{geometry_6.png}
\caption{Load case 6.}
\end{figure}
\clearpage

\begin{figure}[htbp]
\centering
\includegraphics[width=\textwidth]{geometry_7.png}
\caption{Load case 7.}
\end{figure}

\begin{figure}[htbp]
\centering
\includegraphics[width=\textwidth]{geometry_8.png}
\caption{Load case 8.}
\end{figure}

\begin{figure}[htbp]
\centering
\includegraphics[width=\textwidth]{geometry_9.png}
\caption{Load case 9.}
\end{figure}
\clearpage

\begin{figure}[htbp]
\centering
\includegraphics[width=\textwidth]{geometry_10.png}
\caption{Load case 10.}
\end{figure}


\subsubsection{The centre-load cases, which neither code reproduces}

Cases 2, 3, 7 and 8 put \SI{9.7}{} or \SI{25.2}{kg} on the mid-span leading
edge --- a node with \emph{no bridle attached}, carrying up to twice the
kite's own weight. Measurement has the span slightly \emph{widening} against
bare gravity (\SI{5.04}{} against \SI{4.87}{m}): the kite tilts forward and
the tips fold out.

Neither code reproduces that. \texttt{kite\_fem} overstates the span by
34--48\%, and its author is explicit about the cause: ``by applying such a
local load directly on the inflatable beam, the cross-section deforms and the
stiffness of the tube reduces \dots\ the simplification made by modelling such
a complex structure as a beam element cannot capture this'' (thesis \S7.1).
Billow fails differently and more visibly: it converges --- to residuals of
order \SI{1e-8}{N} --- onto a \emph{folded} branch, closing the span to under
\SI{1}{m}. Seeding from the gravity equilibrium and ramping the load in does
not recover the physical branch. The energy is non-convex, and this is a
second local minimum, not a numerical failure: a converged wrong answer.

These four cases are marked $\dagger$ in Table~\ref{tab:hang-results} and reported separately
rather than averaged in, because a mean over them measures how each code
fails rather than how well it works.

\clearpage
\subsection{What this case actually tests}
\label{sec:hang-scope}

The comparison was built to isolate the canopy formulation, since that is the
one element type that differs between the two codes. It does not do that. A
sensitivity sweep on case~1:

\begin{center}
\begin{tabular}{lr}
\toprule
change & span (m) \\
\midrule
baseline & 7.19 \\
canopy $Et$ \SI{5000}{}$\rightarrow$\SI{1000}{N/m} & 7.15 \\
canopy deleted entirely ($Et=\SI{1}{N/m}$) & 7.36 \\
tube moment--curvature curve scaled by $0.3$ & \textbf{3.73} \\
\bottomrule
\end{tabular}
\end{center}

Deleting the canopy outright moves the answer less than a quarter as far as a
modest change in tube bending. \textbf{This is a tube-bending test}, and any
canopy conclusion read off it would be unsupported.

\paragraph{Collapse.} None of the three codes models post-collapse behaviour.
Breukels' source data characterises it --- the measured moment falls past
collapse --- but \texttt{kite\_fem} computes a collapse flag and never acts on
it (a literal \texttt{TODO: verify collapse} in \texttt{BeamElement.py}), and
Billow saturates at $M_\text{max}$ instead of dropping. A law that does follow
collapse, fixed uniquely by requiring it to match both the fitted initial
stiffness and the fitted collapse curvature,
\begin{equation}
W(\kappa) = EI_0\,\kappa_c^{2}
\left[1-\left(1+\tfrac{\kappa}{\kappa_c}\right)e^{-\kappa/\kappa_c}\right],
\qquad
M(\kappa) = EI_0\,\kappa\,e^{1-\kappa/\kappa_c},
\end{equation}
\emph{converges without difficulty} (56--77 iterations, $\sim$\SI{3}{s},
residual \num{3e-9}) and changes the answer by \SI{20}{mm}. It changes nothing
here because Billow's beams reach only 31\% of the collapse curvature and none
of the 98 is flagged collapsed. Note also that collapse arrives at
$0.7\,\kappa_0$, i.e.\ at only $\sim$52\% of $M_\text{max}$, so the saturating
plateau is already too generous well before the cliff. It matters in general;
it does not matter in this dataset.

\subsection{Cost, convergence, and what the \texttt{kite\_fem} numbers are}

\begin{table}[htbp]
\centering
\caption{Solver cost. $^{\ast}$ marks cases that never reached
\texttt{kite\_fem}'s own \num{0.01} tolerance. \textbf{The two time columns were
measured on different machines} and are not a controlled benchmark --- see the
caveats below.}
\begin{tabular}{r rr rrr}
\toprule
LC & \multicolumn{2}{c}{\texttt{kite\_fem}} & \multicolumn{3}{c}{Billow} \\
\cmidrule(lr){2-3}\cmidrule(lr){4-6}
 & $t$ (s) & final tol. & $t$ (s) & iterations & $\|r\|$ (N) \\
\midrule
1 & 546 & 0.010 & 7 & 89 & 2e-07 \\
2 & 671 & 0.010 & 15 & 200 & 9e-08 \\
3 & 465 & 0.229$^{\ast}$ & 19 & 228 & 4e-07 \\
4 & 443 & 0.010 & 16 & 212 & 1e-07 \\
5 & 619 & 0.012$^{\ast}$ & 18 & 240 & 4e-07 \\
6 & 576 & 0.010 & 6 & 86 & 4e-09 \\
7 & 621 & 0.010 & 15 & 209 & 4e-09 \\
8 & 468 & 0.076$^{\ast}$ & 17 & 235 & 6e-09 \\
9 & 483 & 0.010 & 15 & 187 & 1e-07 \\
10 & 626 & 0.011 & 16 & 227 & 4e-07 \\
\midrule
mean & 552 & & 14 & 191 & \\
\bottomrule
\end{tabular}
\end{table}

Billow reaches force balance to better than \SI{5e-7}{N} on every case, against
a target of \SI{1e-4}{N}, in a mean of \SI{14}{s}. \texttt{kite\_fem}'s recorded
times average \SI{552}{s} and three of its ten cases stop above its own
tolerance (\num{0.229}, \num{0.076}, and a marginal \num{0.0117}).

Robustness required one addition to the solver. On slack-dominated
configurations a near-zero-curvature search direction is answered with a step of
order \num{1e4}, the objective overflows, and the line search collapses before
restoration can help. A move limit --- a trust region expressed as bounds IPOPT
already enforces --- cures this and reproduces the direct solve to the digit.
One caveat came with it: once bounds are active, IPOPT reports success for the
\emph{boxed} problem while the iterate sits on the boundary, so the acceptance
test must key off the force residual and never off the solver's own verdict.

\subsubsection{Caveats on every \texttt{kite\_fem} number in this section}
\label{sec:hang-caveats}

These matter enough to state plainly, because the comparison is the weakest part
of this work and the reader should discount it accordingly.

\paragraph{It was never reproduced.} Every \texttt{kite\_fem} figure here comes
from its \emph{committed} \texttt{model\_results.csv} and saved \texttt{.npz}
states, not from a run performed for this study. One attempt was made; it
produced no output in over ninety minutes and was abandoned. So nothing here
distinguishes ``\texttt{kite\_fem} as published'' from ``\texttt{kite\_fem} as it
runs today''.

What \emph{was} verified is narrower but load-bearing: our length extractor
reproduces its published table from its own saved states to \SI{1e-9}{m} on all
ten cases (a committed test). The rescoring is therefore sound even though the
solve is not reproduced.

\paragraph{Its solver settings are tuning choices we did not replicate.} The
published run uses an \texttt{I\_stiffness} of 15, a tolerance of \num{0.01}, a
step limit, a relaxation schedule, and an outer loop that calls
\texttt{adapt\_stiffnesses} --- progressively stiffening any element above 1\%
strain, up to four times. None of that was re-derived or assessed here. A
different setting could move its answers by an unknown amount.

\paragraph{It is not solving the model it states.} Its converged beams sit at
about $0.23\times$ its own constitutive law, with negative bending stiffness in
three of the four cases inspected, and its anchor node belongs to no element so
its equilibrium is held by \texttt{I\_stiffness} rather than by a load path ---
its own solutions drift \SIrange{0.29}{0.34}{m} bodily. Comparing against it is
therefore comparing against a numerical state, not a clean reference solution.

\paragraph{``\texttt{kite\_fem}'' is not one thing.} The results used here are
from \texttt{main} at \texttt{630c8c0}. AWETrim's own environment installs the
branch \texttt{mass\_proportional\_regularisation} at \texttt{6edff17} --- named
for the very regularisation implicated above. Any statement in this section about
\texttt{kite\_fem} refers to the former.

\paragraph{The bridle is not matched in the headline comparison.}
\texttt{kite\_fem} ran the bridle exactly as stored. Our primary results apply
the Table~B.1 corrections. Scored on the stored bridle for a like-for-like
comparison, Billow gives 15.1\% against \texttt{kite\_fem}'s 12.4\%; on the
corrected bridle it gives 15.7\%. The difference the bridle makes overall is
therefore \SI{0.6}{} percentage points --- small, but the comparison is not
bridle-matched and should not be read as though it were.

\paragraph{The metric is ours, not theirs.} Its published
\texttt{shapecorrelation} is not used, for the reason given in Table~1's
caption. Numbers here will therefore not match those reported upstream.

\paragraph{The timings are not a benchmark.} Its recorded times are from an
HP~Zbook Power~G7; ours are from a different machine, running Python against a
compiled CasADi/IPOPT stack while \texttt{kite\_fem} runs Cython
\texttt{pyfe3d} elements. Read the ratio as an order of magnitude at most, and
note that the two codes are not even solving to the same tolerance --- a factor
of \num{1e3} separates the residuals.

\subsection{What stiffness the measurement implies}
\label{sec:hang-lambda}

Rather than fit bridle lengths, the whole tube moment--curvature curve was
scaled by a single factor $\lambda$ --- ``\texttt{moment\_max}'',
``\texttt{bending\_stiffness}'' and ``\texttt{torsion\_c1}'' together, so
$\kappa_0 = M_\text{max}/EI_0$ is unchanged and the curve keeps its shape while
only its magnitude moves. The two bare-gravity cases, the only two that
genuinely load the tubes, then say what stiffness the measurement wants:

\begin{center}
\begin{tabular}{rrrrr}
\toprule
$\lambda$ & \multicolumn{2}{c}{case 1, \SI{0.15}{bar}} & \multicolumn{2}{c}{case 6, \SI{0.25}{bar}} \\
\cmidrule(lr){2-3}\cmidrule(lr){4-5}
 & span (m) & error (\%) & span (m) & error (\%) \\
\midrule
1.00 & 7.276 & 25.9 & 7.468 & 19.5 \\
0.55 & 6.156 & 14.6 & 6.589 & 11.5 \\
0.45 & 5.569 &  9.3 & 6.120 &  7.7 \\
\textbf{0.35} & \textbf{4.674} & \textbf{8.3} & \textbf{5.376} & \textbf{5.5} \\
0.25 & 3.376 & 18.6 & 4.141 & 14.5 \\
\midrule
\multicolumn{1}{r}{measured} & 4.867 & & 5.453 & \\
\bottomrule
\end{tabular}
\end{center}

\noindent
Both cases imply $\lambda \approx 0.36$, and --- this is the part that matters
--- they imply the \emph{same} factor at both pressures. A pressure-dependent
error in the fit would not do that; a constant offset would. It is also further
evidence against the low-pressure extrapolation of $M_\text{max}$ discussed
above, which would have required more softening at \SI{0.15}{bar} than at
\SI{0.25}{bar}.

Run across all ten cases, a single $\lambda = 0.36$ gives a mean relative error
of 8.0\% against \texttt{kite\_fem}'s 12.4\%, and 6.2\% against 11.5\% on the two
cases that test the tubes. It also repairs the centre-load cases, which the
unscaled model got worst: case~2 falls from 27.6\% to 9.6\%, case~7 from 19.7\%
to 7.5\%. The tip-load cases barely move (6.5\%$\to$6.6\%, 5.2\%$\to$5.2\%),
exactly as the previous section predicts --- they are bridle-determined and
carry no information about the tube law.

\paragraph{This is an identification, not a fit.} One scalar, applied uniformly
to all 98 tube elements, calibrated on two cases and improving eight. It is
reported because a consistent, pressure-independent factor is evidence about the
fit rather than a free parameter absorbing error --- but it is not adopted as a
model change. Two things must be understood first.

\emph{Where does $0.36$ come from?} Inverting the fitted $EI_0$ through
$EI = E t \pi r^3$ gives an implied $E t$ of \SIrange{133}{155}{kN/m} for the
leading-edge tubes, nearly constant across $d = \SIrange{0.113}{0.202}{m}$ ---
so the fit's radius dependence is physically consistent, and the implied fabric
(\SIrange{0.5}{1.2}{GPa} at \SIrange{0.15}{0.25}{mm}) is plausible to soft for
polyester. The fit is therefore not obviously twice too stiff on material
grounds, which makes a calibration or usage mismatch the more likely
explanation.

Two candidates are worth checking against Breukels' thesis directly. First, his
tubes are discretised as \emph{rigid cylindrical segments joined by torque
springs}, and the $C_1$--$C_8$ coefficients were fitted so that such a chain
reproduces the measured beam --- a discretisation-dependent calibration, not an
intrinsic moment--curvature law, and re-using it as a continuum law could carry
a systematic factor. Second, his nomenclature defines an \emph{effective}
internal pressure $p_\text{eff}$ distinct from the gauge pressure; if the fits
are calibrated on $p_\text{eff}$, feeding \SI{0.15}{bar} gauge straight in is
wrong at exactly this kind of magnitude.

\subsection{How weakly the published data constrains a model}

The twelve published quantities are less information than they look. Every
left/right pair is identical to $0.0\times10^{0}$ --- \texttt{La}$\equiv$%
\texttt{Li}, \texttt{LcsTL}$\equiv$\texttt{LcsTR} --- so the data is exactly
mirrored and only six numbers are independent. The thesis itself scores just
three: the trailing-edge length (the sum of \texttt{La}--\texttt{Li}), the span,
and the tip-to-centre distance.

More limiting is what is absent. There is no height, no sag depth, and no
chordwise measure of any kind: every published quantity lies in the spanwise /
trailing-edge geometry. A model can match all three scored lengths and still
have the wrong shape --- the wrong camber, the wrong chordwise twist, the wrong
depth of sag --- with nothing in the dataset able to detect it. That is a real
ceiling on what any validation against this file can claim, including this one.

The information exists. The measurement was stereoscopic, giving full
three-dimensional marker positions, and the thesis plots measured geometry
against model geometry in top, side and front projection for all ten load cases
(its Figure~7.2). Only the twelve scalars were published.

\paragraph{The single most valuable thing to obtain} is therefore the 3D marker
coordinates behind Figure~7.2. They would turn this from a comparison of three
effective scalars into a comparison of shapes, and would immediately settle
questions the present data cannot touch --- among them whether Billow's wing is
uniformly too open or wrong in a particular region, and whether the centre-load
cases fail in the way the folded branch suggests.

\subsection{The proper test: 3D shape against the markers}
\label{sec:hang-shape}

Everything above is scored on twelve numbers that reduce to six independent
ones. Haanen's marker clouds --- 62 to 76 three-dimensional points per load case
on the leading edge, the eight struts and the canopy --- allow the question that
data cannot answer: is the \emph{shape} right?

\paragraph{Method.} Markers are in the stereo rig's frame, so a \emph{rigid}
transform is needed --- rotation and translation only, never scale, since scale
is what is being tested. Correspondence comes from structure rather than
fitting: the measured struts are ordered across the span with their first marker
at the leading edge, so their roots pair with the model's eight strut
leading-edge nodes, giving eight anchors for a Kabsch fit. Both spanwise
orderings are tried, because the kite is near-symmetric and the marker numbering
direction is undocumented. The fit is then refined by ICP against the model
curves, so no marker is assumed to sit on a node. Reported errors are
point-to-\emph{surface}: the markers are stuck to the outside of the tubes while
the model polyline is the tube axis, and ignoring that radius --- 55 to
\SI{100}{mm} here --- would put a systematic offset into every case.

\begin{table}[htbp]
\centering
\caption{Point-to-surface RMS between the solved shape and the photogrammetry
markers. ``As built'' is the undeformed geometry, included as the null model:
anything not beating it has explained nothing.}
\begin{tabular}{rl rrr}
\toprule
LC & load & \multicolumn{2}{c}{Billow (mm)} & as built (mm) \\
\cmidrule(lr){3-4}
 & & EI x1.0 & EI x0.36 & \\
\midrule
1 & gravity & 480 & 100 & 695 \\
2 & centre 9.7 kg & 501 & 152 & 654 \\
3 & centre 25.2 kg & 507 & 171 & 639 \\
4 & tip 2 kg & 144 & 136 & 141 \\
5 & tip 5 kg & 101 & 92 & 102 \\
6 & gravity & 387 & 101 & 554 \\
7 & centre 9.7 kg & 399 & 151 & 525 \\
8 & centre 25.2 kg & 409 & 145 & 516 \\
9 & tip 2 kg & 119 & 111 & 114 \\
10 & tip 5 kg & 102 & 94 & 101 \\
\midrule
\multicolumn{2}{r}{mean, gravity only} & \textbf{434} & \textbf{101} & 625 \\
\multicolumn{2}{r}{mean, centre load} & \textbf{454} & \textbf{155} & 583 \\
\multicolumn{2}{r}{mean, tip load} & \textbf{117} & \textbf{108} & 115 \\
\multicolumn{2}{r}{mean, all ten} & \textbf{315} & \textbf{125} & 404 \\
\bottomrule
\end{tabular}
\end{table}

\paragraph{The scaled model reaches the measurement floor.} On the two cases
that load the tubes, the error falls from \SI{480}{} to \SI{100}{mm} and from
\SI{387}{} to \SI{101}{mm} --- onto the same $\sim$\SI{100}{mm} the tip cases
sit at, which is the marker placement and alignment floor. Across all ten the
mean falls from \SI{315}{} to \SI{125}{mm}, against \SI{404}{mm} for doing
nothing. That the factor identified from two scalar spans also fixes the full
three-dimensional shape, including the centre-load cases it was not fitted on,
is the strongest evidence in this section that it is a real property of the tube
law rather than a number absorbing error.

\paragraph{The tip-load cases are confirmed uninformative.} Their errors are
\SI{141}{}, \SI{102}{}, \SI{114}{}, \SI{101}{mm} for the \emph{undeformed}
geometry against \SI{144}{}, \SI{101}{}, \SI{119}{}, \SI{102}{mm} for the
unscaled model and \SI{136}{}, \SI{92}{}, \SI{111}{}, \SI{94}{mm} for the scaled
one --- all within about \SI{10}{mm} of each other. Doing nothing scores as well
as either model. This is the same conclusion the bridle-tension argument reached
in Section~\ref{sec:hang-scope}, now established directly in three dimensions rather than inferred,
and it is why the six-case and ten-case means elsewhere in this section should be
read with care.

\subsection{Summary}

\begin{itemize}
\item The distributed reference model is ill posed --- its anchor belongs to no
      element --- and its published results are held up by a numerical
      regularisation. Any future use of this dataset has to supply a suspension.
\item \texttt{kite\_fem}'s beams converge to $\sim0.23\times$ its own
      constitutive law, with negative bending stiffness in several cases. Its
      closer agreement with the measurement is not evidence of better physics.
\item The dataset is a tube-bending test --- and only two of its ten cases even
      are. It cannot decide a canopy model, which was the original reason for
      building the comparison. The 3D markers confirm this directly: on the
      tip-load cases the undeformed geometry scores as well as either model.
\item Against the 3D markers, $\lambda = 0.36$ takes the mean point-to-surface
      error from \SI{315}{} to \SI{125}{mm} (null model \SI{404}{mm}), and
      brings the two tube-loaded cases onto the $\sim$\SI{100}{mm} measurement
      floor --- having been identified on two scalar spans, not on the shape.
\item Post-collapse behaviour is unmodelled in all three codes and does not
      matter here or in coupled use, because the tubes never leave their initial
      slope.
\item The residual disagreement lives in $EI_0$. A single pressure-independent
      factor $\lambda \approx 0.36$, identified on the two gravity cases,
      improves eight of the ten and beats \texttt{kite\_fem} overall
      (8.0\% against 12.4\%). It is reported as an identification, not
      adopted: two specific mismatches in how the Breukels fit is being
      re-used --- its torque-spring discretisation and its effective
      pressure --- should be checked against the source first.
\end{itemize}


\section{Demonstration cases}

\subsection{Wrinkling: fabric in a frame smaller than itself}

\begin{figure}[H]
\centering
\includegraphics[width=\textwidth]{wrinkling.png}
\caption{A panel clamped into a frame 6\% smaller than itself, then pushed out
of plane. The relaxed law converges in 31 iterations to a smooth billow; the
unrelaxed one takes 211 and buckles into mesh-scale crumple at less than
two-thirds the deflection.}
\label{fig:wrinkling}
\end{figure}

This is the clearest illustration of why the relaxed energy is not optional.
The unrelaxed panel stores $40.6$\,J against the relaxed panel's $0.116$\,J,
almost all of it fighting compression that fabric cannot carry.

\subsection{Large-deflection cantilever}

\begin{figure}[H]
\centering
\includegraphics[width=\textwidth]{cantilever.png}
\caption{Deflected shapes under a tip load swept from 50 to 900\,N, and tip
deflection against linear theory. The beam visibly shortens in $x$ as it bends
--- the geometric nonlinearity linear theory omits. At the largest load, linear
theory overshoots by 51\%.}
\label{fig:cantilever}
\end{figure}

Median 9 IPOPT iterations per load step with warm starting.

\subsection{Batten-stiffened sail}

\begin{figure}[H]
\centering
\includegraphics[width=\textwidth]{sail.png}
\caption{Beams, membranes, cables and a trailing-edge hem in one energy: a
spar carrying bending and torsion, a wrinkling canopy hanging off it, three
bridles to a fixed anchor. 393 DOF, 54 iterations, 1.4\,s, residual
$1.0\times10^{-9}$\,N.}
\label{fig:sail}
\end{figure}

Note that only 13 of the nodes carry a material frame; the canopy nodes stay at
three degrees of freedom.

\subsection{Scaling}

\begin{figure}[H]
\centering
\includegraphics[width=\textwidth]{scaling.png}
\caption{Build and solve time against problem size, and iteration count.
Because the graph is built once per element \emph{type}, build cost tracks the
solve rather than the element count, and the iteration count is nearly mesh
independent across a 37$\times$ range in degrees of freedom.}
\label{fig:scaling}
\end{figure}

%======================================================================
\section{Comparison with \texorpdfstring{\code{kite\_fem}}{kite\_fem}}

\code{kite\_fem} (with \code{pyfe3d}) is the existing FEM path in the AWETrim
tree. This section states what differs, what was measured, and where the
comparison is and is not fair.

\subsection{Capability}

\begin{table}[H]
\centering
\small
\begin{tabular}{p{4.0cm}p{5.1cm}p{5.1cm}}
\toprule
& \code{kite\_fem} / \code{pyfe3d} & \textbf{Billow}\\
\midrule
Elements & springs, pulleys, inflatable Timoshenko beams
         & cables, pulleys, geometrically exact Timoshenko beams,
           inflatable tubes, wrinkling membranes\\
Canopy fabric & noncompressive spring net (no membrane element exists)
              & CST membrane, Pipkin relaxed energy\\
Degrees of freedom & 6 per node everywhere, masked via \code{bc}
                   & 3 per node, $+3$ only on beam nodes\\
Formulation & Newton--Raphson on $\mathbf{f}_e-\mathbf{f}_i$
            & minimum potential energy (IPOPT)\\
Assembly & scipy sparse COO per element & one mapped kernel per element type\\
Regularisation & \code{I\_stiffness=25}, \code{step\_limit=0.2},
                 \code{relax=0.5}, optional pseudo-transient
               & IPOPT inertia correction; membrane slack blend\\
Warm start & none (\code{solve} resets to undeformed) & state and duals\\
Default tolerance & $10^{-2}$\,N & $10^{-8}$ IPOPT, accepted at $10^{-6}$\,N\\
Inflatable tube & secant $EI$/$GJ$ plus collapse flag
                & the same fits integrated to $W(\kappa)$\\
Experimental validation & hanging-kite test against measured data & none\\
\bottomrule
\end{tabular}
\end{table}

\subsection{Two traps in making it a fair comparison}

\begin{enumerate}
\item \code{FEM\_structure.solve} defaults to \code{I\_stiffness=25}, an
\textbf{absolute} regularisation in N/m added to every degree of freedom of the
tangent. The benchmark beams here have $EI/L^3 = 10$\,N/m, so the default is
larger than the structure: 2000 iterations with the residual stuck at 20\,N.
With it at zero the same problem converges in 38. It is a kite-scale tuning
constant, not a universal one. All results above use zero.
\item The property matching was verified in the linear limit before any
conclusion was drawn: below about one degree of rotation the two codes agree to
$0.02\%$.
\end{enumerate}

\subsection{Cost and accuracy on a common problem}

Tip-loaded slender cantilever, matched properties, same mesh, one machine, one
run.

\begin{table}[H]
\centering
\small
\caption{Source: \code{run\_cost\_comparison.py}. ``Warm'' is a second solve at
a 5\% different load; \code{kite\_fem} restarts from the undeformed
configuration on every call, so its warm column is a second cold solve.}
\begin{tabular}{rrlrrrrrr}
\toprule
$\alpha$ & El. & Model & DOF & Setup (s) & Cold (s) & Warm (s) & Iter. & ms/iter\\
\midrule
1 & 20 & Billow           & 126 & 0.133 & 0.249 & \textbf{0.081} & 16 & 15.6\\
1 & 20 & \code{kite\_fem} & 126 & \textbf{0.004} & \textbf{0.132} & 0.151 & 39 & \textbf{3.4}\\
1 & 80 & Billow           & 486 & 0.217 & 0.733 & \textbf{0.282} & 17 & 43\\
1 & 80 & \code{kite\_fem} & 486 & \textbf{0.009} & \textbf{0.189} & 0.192 & 39 & \textbf{4.8}\\
6 & 40 & Billow           & 246 & 0.076 & \textbf{0.333} & \textbf{0.062} & \textbf{22} & 15\\
6 & 40 & \code{kite\_fem} & 246 & \textbf{0.005} & 1.384 & 1.526 & 444 & \textbf{3.1}\\
6 & 80 & Billow           & 486 & 0.186 & \textbf{0.669} & \textbf{0.233} & \textbf{23} & 29\\
6 & 80 & \code{kite\_fem} & 486 & \textbf{0.005} & 2.270 & 2.658 & 490 & \textbf{4.8}\\
\bottomrule
\end{tabular}
\end{table}

\code{kite\_fem} has roughly $6\times$ cheaper iterations (compiled \code{pyfe3d}
plus scipy sparse) and $30\times$ faster setup. Billow uses 2--20$\times$ fewer
iterations, and its count is nearly independent of both mesh and load
(16--23 throughout) where \code{kite\_fem}'s climbs from 39 to about 490 as the
load grows. \code{kite\_fem} therefore wins on easy, small, one-shot problems;
Billow wins as soon as the problem is hard, reused, or large. Billow's setup
cost is a one-off amortised across a sweep --- exactly the coupled-loop pattern.

\subsection{The canopy model: spring net versus membrane}

Since \code{kite\_fem} has no membrane element, its canopy is a noncompressive
spring net (chordwise and spanwise links, \emph{no diagonals}; its own example
carries a commented-out block of ``fictional springs constraining the canopy from
collapse''). Billow can build exactly that net from its cable elements, which
makes a genuine like-for-like study possible.

\begin{figure}[H]
\centering
\includegraphics[width=\textwidth]{canopy_model.png}
\caption{Spring net versus CST membrane on identical meshes, both minimised.
Left: both converge under out-of-plane load, to answers 7\% apart. Middle: under
equibiaxial stretch the net is too soft by exactly $1-\nu$. Right: under in-plane
shear the membrane energy scales as $\gamma^{2.0}$ and the net as $\gamma^{4.0}$
--- the net has \emph{no shear stiffness at leading order}.}
\label{fig:canopymodel}
\end{figure}

Measured: biaxial energy ratio $0.728$ against $1-\nu = 0.70$; shear exponents
$2.018$ (membrane) and $3.999$ (net).

\subsubsection{Taut: shape and cost}

\begin{figure}[H]
\centering
\includegraphics[width=\textwidth]{taut_net_vs_membrane.png}
\caption{Taut bay, $24\times18$, both minimised on the same mesh with the same
1425 degrees of freedom, so this is purely the canopy model. The net billows
$9.6\%$ deeper; the best single scale factor is $1.120$, and a $3.31$\,mm
residual shape difference remains after removing it.}
\label{fig:tautshape}
\end{figure}

The local softness ratio is not constant: $1.10$ at mid-chord rising to $1.145$
near the clamped edges, where the fabric carries biaxial tension and shear that
the net cannot represent. The residual map shows the net relatively too deep in
the corners and too shallow in the middle --- it billows into a slightly
different surface, not merely a deeper one.

\textbf{The scale factor cannot be calibrated away}, because it depends on the
stress state: about $1.10$--$1.12$ when taut, but $1.00$ when fully wrinkled
(the slack cases agree within $1\%$, since a deeply slack panel carries
essentially uniaxial chordwise tension where the two models coincide). The
correction needed would change with flight condition.

Cost is not a differentiator. At $24\times18$: net $0.092$\,s setup and
$0.227$\,s solve in 12 iterations; membrane $0.176$\,s and $0.245$\,s in 8. The
membrane's per-iteration cost is higher but it needs fewer iterations, and the
two roughly cancel.

\subsubsection{Robustness}

\begin{table}[H]
\centering
\small
\caption{Peak billow (\% chord) and section RMS against each model's own finest
mesh. The net fails outright at one intermediate slack mesh --- the signature of
its zero-energy shear mode. The membrane never failed.}
\begin{tabular}{lrrrrrr}
\toprule
& \multicolumn{3}{c}{Spring net} & \multicolumn{3}{c}{Membrane}\\
\cmidrule(lr){2-4}\cmidrule(lr){5-7}
Mesh & Camber & RMS & Conv. & Camber & RMS & Conv.\\
\midrule
\multicolumn{7}{l}{\emph{taut}}\\
$4\times3$   & 2.858 & 0.388 & yes & 2.542 & 0.401 & yes\\
$8\times6$   & 3.129 & 0.051 & yes & 2.842 & 0.051 & yes\\
$16\times12$ & 3.117 & 0.010 & yes & 2.842 & 0.010 & yes\\
$24\times18$ & 3.115 & ---   & yes & 2.842 & ---   & yes\\
\multicolumn{7}{l}{\emph{slack 3\%}}\\
$8\times6$   & 7.388 & 0.186 & yes & 7.449 & 0.143 & yes\\
$12\times9$  & 7.240 & 0.091 & yes & 7.294 & 0.097 & yes\\
$16\times12$ & 5.675 & 1.719 & \textbf{no} & 7.412 & 0.031 & yes\\
$24\times18$ & 7.302 & ---   & yes & 7.402 & ---   & yes\\
\bottomrule
\end{tabular}
\end{table}

\subsubsection{Canopy model versus solver, separated}

A single bay isolates the two effects that otherwise get conflated. Rows 2 and 3
solve the \emph{identical} spring net, so their difference is the solver alone.

\begin{table}[H]
\centering
\small
\caption{One $2.0\times1.3$\,m bay at 250\,Pa, all edges fixed.}
\begin{tabular}{llrrrr}
\toprule
Case & Configuration & Solve (s) & Iterations & Camber (\% c) & Converged\\
\midrule
taut $8\times6$ & Billow + membrane   & 0.030 & 7 & 2.842 & yes\\
taut $8\times6$ & Billow + net        & 0.026 & 10 & 3.129 & yes\\
taut $8\times6$ & \code{kite\_fem} + net & 10.04 & 2729 & 3.129 & yes\\
\midrule
taut $12\times9$ & Billow + net        & 0.034 & 9 & 3.083 & yes\\
taut $12\times9$ & \code{kite\_fem} + net & 22.00 & 3001 & 5.283 & \textbf{no}\\
\midrule
slack $8\times6$ & Billow + net        & 0.049 & 23 & 7.388 & yes\\
slack $8\times6$ & \code{kite\_fem} + net & 9.06 & 3001 & 7.363 & \textbf{no}\\
\bottomrule
\end{tabular}
\end{table}

Where both converge they agree exactly ($3.129\%$ both), which is what makes the
\code{kite\_fem} driver used throughout this document trustworthy. On the
identical net, minimum energy is roughly $390\times$ faster than the damped
Newton solve, and remains convergent where the Newton solve does not.

\subsection{Full kite}

A V3-scale structure --- inflated leading-edge tube, inflated struts at every bay
boundary, canopy over nine bays, trailing-edge hem, bridle to a fixed anchor,
uniform 250\,Pa --- was built with the same nodes, tubes and cables in both codes.

\begin{table}[H]
\centering
\small
\begin{tabular}{lrrrrrl}
\toprule
Model & DOF & Setup (s) & Cold (s) & Warm (s) & Iter. & Residual (N)\\
\midrule
\textbf{Billow} & 1893 & 0.61 & \textbf{5.91} & \textbf{0.58} & 73
 & $7.2\times10^{-9}$ (converged)\\
\code{kite\_fem} ($I{=}25$) & 2976 & 0.05 & 76.2 & 72.0 & 2001 (cap)
 & $4.1\times10^{5}$ (diverged)\\
\code{kite\_fem} ($I{=}0$) & 2976 & 0.06 & 108.9 & 68.9 & 2001 (cap)
 & $1.3\times10^{8}$ (diverged)\\
\bottomrule
\end{tabular}
\end{table}

\paragraph{Interpretation, carefully.} No usable \code{kite\_fem} full-kite
number was obtained. Diagonals, canopy curvature, both regularisation settings
and load stepping with rebuilt initial conditions were all tried; all diverged,
with residuals roughly $100\times$ the applied load. \textbf{This is consistent
with a known open issue on the \code{kite\_fem} side} rather than a defect in the
driver used here --- the ASKITE integration commit records ``kite\_fem still
diverging, but not due to the coupling''. The only kite-scale figure that comes
from its own workflow is its hanging-kite validation, at 442--671\,s per load
case. The single-bay comparison above is the trustworthy cost comparison.

%======================================================================
\section{How many membrane triangles?}

Chordwise and spanwise resolution are swept independently, because they are not
the same problem: a canopy bay billows across the chord --- which sets the
section the aerodynamics sees --- while spanwise it varies slowly between struts.

The test case is one bay of the LEI-V3 canopy: chord 2.0\,m, strut spacing
1.3\,m (the V3 spans 8.28\,m over 10 stations), held on all four edges, under
250\,Pa. Both a taut panel and a realistically slack one (reference cut 3\%
oversize) are run.

Convergence is judged on quantities the coupling actually consumes: peak billow
over chord, and the RMS deviation of the mid-span chordwise profile from the
finest mesh. \textbf{Peak local strain is deliberately not a criterion}: in a
wrinkling membrane it does not converge, and nothing downstream reads it.

\begin{figure}[H]
\centering
\includegraphics[width=\textwidth]{mesh_shapes.png}
\caption{One bay under uniform refinement. A $2\times2$ bay is a tent, not a
canopy: the camber is only 2.8\% low but the section is two straight lines.
By $8\times8$ the profile is visually converged.}
\label{fig:meshshapes}
\end{figure}

\begin{figure}[H]
\centering
\includegraphics[width=\textwidth]{mesh_requirement.png}
\caption{Camber and section-shape convergence, chordwise (solid) and spanwise
(dashed), and cost per bay.}
\label{fig:meshreq}
\end{figure}

\begin{table}[H]
\centering
\caption{Candidate meshes against a $24\times24$ reference, slack 3\%.}
\begin{tabular}{lrrrrr}
\toprule
Mesh & Triangles & Camber error & Section RMS (\% c) & Build $+$ solve (s) & Iter.\\
\midrule
$4\times4$   & 32  & 0.88\% & 0.967 & 0.03 & 12\\
$6\times6$   & 72  & 0.61\% & 0.338 & 0.07 & 17\\
\textbf{$8\times6$} & \textbf{96} & \textbf{0.71\%} & \textbf{0.145} & \textbf{0.09} & 19\\
$8\times8$   & 128 & 0.16\% & 0.140 & 0.17 & 15\\
$12\times8$  & 192 & 0.18\% & 0.075 & 0.18 & 21\\
$16\times12$ & 384 & 0.21\% & 0.033 & 0.44 & 22\\
\bottomrule
\end{tabular}
\end{table}

\paragraph{Recommendation.} $\mathbf{8}$ \textbf{chordwise} $\times$
$\mathbf{6}$ \textbf{spanwise per bay}, 96 triangles, about 900 for the whole
nine-bay canopy. Camber error $0.71\%$ and section RMS $0.145\%$ of chord ---
2.9\,mm on a 2\,m chord, comfortably below what an airfoil parameterisation
resolves. Use $12\times8$ for final runs: it halves the shape error for twice
the cost.

\paragraph{Slack costs roughly twice the mesh.} At 3\% oversize the panel is
88\% wrinkled (against 1--4\% taut) and needs about $12\times8$ for the same
accuracy, with 4$\times$ the iterations. Since a real canopy flies slack, the
as-cut slack of the actual sail drives the mesh requirement more than anything
else here and is worth pinning down.

\paragraph{Caveat.} These numbers are for \emph{uniform} pressure. A real
chordwise pressure distribution is peaked near the leading edge and may demand
more chordwise resolution; the study should be repeated once loads come from an
aerodynamic solve.

%======================================================================
\section{Known limitations}

\begin{description}[style=unboxed,leftmargin=1.2em]
\item[The measured validation cannot decide the canopy.]
Section~\ref{sec:hanging} scores Billow against a real kite, but that dataset
is governed by inflatable-tube bending: deleting the canopy entirely moves the
span by \SI{0.17}{m} while a modest change in tube stiffness moves it by
\SI{2.5}{m}. So the claim that a wrinkling membrane is more correct than a
spring net still rests on it being the physically appropriate model, not on
measurement. A test that loads the canopy in its own right --- pressure normal
to a restrained bay, with the deformed surface measured --- is the missing
experiment.

\item[Tube stiffness is unresolved at roughly a factor of three.] The same
section identifies $\lambda \approx 0.36$ on the moment--curvature curve, which
reproduces the measured three-dimensional shape but is \emph{not} adopted: the
likelier cause is a mismatch in how the Breukels fit is re-used rather than a
wrong fit. Until that is settled, absolute deflections from any Billow model
using \code{InflatableTubeLaw} --- including coupled aero-structural runs, whose
tubes sit on the same initial slope --- carry that uncertainty.

\item[Applied moments.] Need load stepping, and are exact only about a fixed
axis; a dead three-dimensional moment is non-conservative and has no potential.
Nodal forces are unaffected.

\item[Post-collapse tubes.] Reported, not modelled. The fits cover bending and
torsion only; axial and shear stiffness must be supplied from tube geometry.

\item[Dead loads.] Loads are frozen per solve, which matches a staggered
coupling loop. A true follower pressure would not be conservative and could not
be posed as a potential --- worth remembering before anyone tries to fold
aerodynamics into the same minimisation.

\item[No contact.] Contact and self-contact are not modelled.

\item[Setup cost.] About $0.2$\,s of graph construction for a 486-DOF beam is
disproportionate against \code{kite\_fem}'s $0.005$\,s, even though it is a
one-off and stops mattering at membrane scale.

\item[Threaded assembly.] \code{PotentialEnergy(parallelization=...)} exposes
CasADi threaded maps but is untested at scale; serial is the default.

\item[Not coupled.] There is no aerodynamic coupling adapter yet.
\end{description}

%======================================================================
\section{Reproducing everything in this document}

\begin{table}[H]
\centering
\small
\begin{tabular}{ll}
\toprule
Figure / table & Command\\
\midrule
Figs.~\ref{fig:wrinkling}--\ref{fig:scaling}, \ref{fig:canopymodel}
 & \code{python scripts/structural/run\_demo\_cases.py}\\
Figs.~\ref{fig:elastica}--\ref{fig:inflatable}
 & \code{python scripts/structural/run\_validation\_benchmarks.py}\\
Cantilever cost table & \code{python scripts/structural/run\_cost\_comparison.py}\\
Figs.~\ref{fig:meshshapes}, \ref{fig:meshreq}
 & \code{python scripts/structural/run\_mesh\_requirement.py}\\
Single-bay comparison & \code{python scripts/structural/run\_panel\_cost.py}\\
Full-kite table & \code{python scripts/structural/run\_full\_kite\_cost.py}\\
Hanging-kite results, Tables & \code{python scripts/structural/run\_hanging\_validation.py}\\
Stiffness scan & \code{python scripts/structural/run\_hanging\_stiffness\_scan.py}\\
3D marker comparison & \code{python scripts/structural/run\_hanging\_shape\_comparison.py}\\
Hanging figures & \code{python scripts/structural/plot\_hanging\_validation.py}\\
Generated \code{.tex} tables & \code{python scripts/structural/make\_hanging\_report.py}\\
Test suite & \code{pytest tests/structural/}\\
\bottomrule
\end{tabular}
\end{table}

\subsection{Package layout}

\begin{table}[H]
\centering
\small
\begin{tabular}{ll}
\toprule
\code{rotations.py} & $SO(3)$ kernels over an \code{xp} namespace (NumPy or CasADi)\\
\code{model.py} & \code{DofLayout}, \code{StructuralState}, \code{StructuralModel}\\
\code{energy.py} & \code{PotentialEnergy}: mapped assembly, gradient, tangent stiffness\\
\code{solver.py} & \code{MinimumEnergySolver}, \code{StructuralSolution}\\
\code{elements/base.py} & \code{ElementKernel} protocol, \code{ElementSet}\\
\code{elements/cable.py} & \code{CableKernel}, \code{PulleyKernel}\\
\code{elements/beam.py} & \code{TimoshenkoBeamKernel}, \code{BeamSection}\\
\code{elements/inflatable.py} & \code{InflatableTubeLaw}, \code{InflatableBeamKernel}\\
\code{elements/membrane.py} & \code{MembraneKernel}, \code{membrane\_regimes}\\
\bottomrule
\end{tabular}
\end{table}

\subsection{Minimal example}

\begin{quote}\small\ttfamily
\noindent
import numpy as np\\
from awetrim.structural import MinimumEnergySolver, StructuralModel\\
from awetrim.structural.elements import build\_cable\_elements\\[4pt]
nodes = np.array([[0., 0, 0], [1., 0, 0]])\\
cables = build\_cable\_elements([[0, 1]], rest\_lengths=[1.0], stiffness=[1e4])\\
model = StructuralModel(nodes, [cables], fixed\_translation\_nodes=[0])\\
solution = MinimumEnergySolver(model).solve(np.array([[0., 0, 0], [10., 0, 0]]))
\end{quote}

\subsection{Invariants to preserve}

\begin{itemize}
\item \textbf{Isolation.} NumPy and CasADi only.
\item \textbf{One mapped kernel per element type.} Never reintroduce a Python
loop over elements that builds \code{MX}.
\item \textbf{Single source for every formula.} \code{rotations.py} holds the
$SO(3)$ maps against an \code{xp} namespace so NumPy and CasADi share one
implementation; \code{beam\_strains} and \code{green\_strain} are called by both
the kernels and the NumPy setup and diagnostic paths.
\item \textbf{No CasADi across the module boundary.} Symbolics are created inside
\code{energy.py} and stay there.
\item \textbf{Kernels are pure.}
\item \textbf{\code{converged} is the physical verdict}, with IPOPT's own kept
alongside.
\end{itemize}

%======================================================================
\section*{References}
\addcontentsline{toc}{section}{References}

\begin{enumerate}[label={[\arabic*]}]
\item J.C. Simo and L. Vu-Quoc, ``A three-dimensional finite-strain rod model.
Part II: computational aspects'', \emph{Comput. Methods Appl. Mech. Engrg.}
\textbf{58} (1986) 79--116.
\item A. Ibrahimbegovi\'c, ``On the choice of finite rotation parameters'',
\emph{Comput. Methods Appl. Mech. Engrg.} \textbf{149} (1997) 49--71.
\item K.-J. Bathe and S. Bolourchi, ``Large displacement analysis of
three-dimensional beam structures'', \emph{Int. J. Numer. Meth. Engng.}
\textbf{14} (1979) 961--986.
\item A.C. Pipkin, ``The relaxed energy density for isotropic elastic
membranes'', \emph{IMA J. Appl. Math.} \textbf{36} (1986) 85--99.
\item D.G. Roddeman, J. Drukker, C.W.J. Oomens and J.D. Janssen, ``The wrinkling
of thin membranes'', \emph{J. Appl. Mech.} \textbf{54} (1987) 884--892.
\item K.E. Bisshopp and D.C. Drucker, ``Large deflection of cantilever beams'',
\emph{Q. Appl. Math.} \textbf{3} (1945) 272--275.
\item Y. Luo, ``An efficient 3D Timoshenko beam element with consistent shape
functions'', \emph{Adv. Theor. Appl. Mech.} \textbf{1}(3) (2008) 95--106.
\item O. Cayon, M. Gaunaa and R. Schmehl, ``Fast aero-structural model of a
leading-edge inflatable kite'', \emph{Energies} \textbf{16} (2023) 3061.
\item ASKITE, \url{https://github.com/awegroup/ASKITE};
\code{kite\_fem}, \url{https://github.com/awegroup/kite_fem}.
\end{enumerate}

\end{document}
